\section*{Bab \ref{ch:g4:division} --- Pembagian}

\begin{solution}{exo:g4:division:1}
$96 \div 4$: antara $4 \times 20 = 80$ dan $4 \times 30 = 120$, jadi
\omterm{def:g3:division:remainder}{hasil baginya} punya dua angka; pembagian bersusun memberi $24$,
\omterm{def:g3:division:remainder}{sisa} $0$.

$87 \div 5$: antara $5 \times 10 = 50$ dan $5 \times 20 = 100$;
\omterm{def:g3:division:remainder}{hasil bagi} $17$, \omterm{def:g3:division:remainder}{sisa} $2$ ($87 = 5 \times 17 + 2$).
\end{solution}

\begin{solution}{exo:g4:division:2}
$672 \div 4 = 168$, \omterm{def:g3:division:remainder}{sisa} $0$ (periksa: $4 \times 168 = 672$).

$925 \div 7 = 132$, \omterm{def:g3:division:remainder}{sisa} $1$ (periksa:
$7 \times 132 + 1 = 925$).

$804 \div 6 = 134$, \omterm{def:g3:division:remainder}{sisa} $0$ (periksa: $6 \times 134 = 804$).
\end{solution}

\begin{solution}{exo:g4:division:3}
$816 \div 4 = 204$ (langkah puluhannya $1 \div 4$: \omterm{def:g1:counting:zero}{nol} kali ---
tulislah angka $0$).

$420 \div 7 = 60$.

$2\,512 \div 5 = 502$, \omterm{def:g3:division:remainder}{sisa} $2$ (periksa:
$5 \times 502 + 2 = 2\,512$).
\end{solution}

\begin{solution}{exo:g4:division:4}
$85 = 9 \times 9 + 4$ (\omterm{def:g3:division:remainder}{hasil bagi} $9$, \omterm{def:g3:division:remainder}{sisa} $4$).

$1\,000 = 3 \times 333 + 1$.
\end{solution}

\begin{solution}{exo:g4:division:5}
$156 \div 6 = 26$: dua puluh enam kotak terisi penuh, tidak ada yang tersisa.
\end{solution}

\begin{solution}{exo:g4:division:6}
$250 = 8 \times 31 + 2$: tiga puluh satu potongan, dan $2$ cm pita
yang tersisa.
\end{solution}

\begin{solution}{exo:g4:division:7}
$375 = 52 \times 7 + 11$: tujuh bus mengangkut $364$ murid, $11$
masih tersisa --- inilah keadaan ``ditambah satu'': $8$ bus diperlukan.
\end{solution}

\begin{solution}{exo:g4:division:8}
$234 \div 3 = 78$: masing-masing mendapat $78$ kelereng. Periksa:
$3 \times 78 = 234$.
\end{solution}

\begin{solution}{exo:g4:division:9}
Bilangan yang dibagi $= 7 \times 58 + 3 = 406 + 3 = 409$.
\end{solution}

\begin{solution}{exo:g4:division:10}
Rak: $438 = 9 \times 48 + 6$: $48$ rak penuh dan $6$ buku lagi ---
jadi $49$ rak diperlukan. Lemari:
$49 = 6 \times 8 + 1$: delapan lemari penuh dan satu rak tambahan
--- jadi $9$ lemari harus dibeli.
\end{solution}

\begin{solution}{exo:g4:division:11}
\omterm{def:g3:division:remainder}{Hasil bagi} $=$ \omterm{def:g3:division:remainder}{sisa} $= r$ dengan $r < 8$, jadi $r$ adalah $0, 1, \dots,
7$ dan bilangan yang dibagi adalah $8r + r = 9r$: bilangan yang
mungkin adalah $0$, $9$, $18$, $27$, $36$, $45$, $54$, $63$.

\end{solution}
